\section{文档概述}

\subsection{报告目的}

本报告说明「寻迹」校园失物招领平台的测试计划、测试用例、缺陷跟踪和质量保证方法。测试工作围绕用户注册登录、校园身份认证、失物与招领发布、搜索匹配、认领交接、积分通知和后台治理展开，用统一的测试准则检查功能行为、权限边界、状态一致性和代码质量。

报告中的测试范围和判定口径与需求、接口、枚举及状态模型保持一致。每个核心场景均给出前置条件、操作步骤和预期结果；缺陷从发现、分级、修复到回归形成闭环；质量保证活动贯穿需求评审、设计评审、编码检查、自动化测试和版本验收。

\subsection{测试对象与范围}

\begin{table}[H]
  \centering
  \caption{测试对象与关注点}
  \begin{tabularx}{\textwidth}{L{2.7cm}L{3.1cm}Y}
    \toprule
    测试对象 & 主要组成 & 重点检查内容 \\
    \midrule
    用户端 & Vue 3、Pinia、Element Plus & 注册登录、发布搜索、匹配认领、个人中心、异常提示与页面状态 \\
    管理端 & Vue 3、管理路由与审核页面 & 管理员鉴权、认证与内容审核、全状态认领审计、举报处置和统计展示 \\
    主后端 & FastAPI、SQLAlchemy、Alembic & 接口契约、业务规则、事务状态、权限、对象访问与异步任务 \\
    AI 服务 & 分类、敏感检测、匹配评分、回答语义校验 & 请求鉴权、输入校验、结果范围、同义表达、异常处理和规则降级 \\
    工程配置 & CI、Compose、环境模板 & 静态检查、自动化测试、生产构建、镜像构建与配置一致性 \\
    \bottomrule
  \end{tabularx}
\end{table}

不属于平台业务范围的短信运营商能力、校园统一身份系统和第三方模型内部算法，不作为本次测试对象；平台对这些外部能力的调用参数、超时处理和异常返回仍纳入接口测试。

\section{测试计划}

\subsection{测试目标}

\begin{enumerate}
  \item \textbf{功能正确性：}确认主要角色能够按规定流程完成操作，输入校验、状态迁移和结果反馈符合需求。
  \item \textbf{权限与隐私：}确认未登录访问、普通用户越权、管理员权限和敏感图片访问均由服务端约束。
  \item \textbf{数据一致性：}确认匹配、认领、交接、积分和通知在重复请求或状态冲突下保持一致。
  \item \textbf{性能与容量：}按照统一负载、数据规模和 P95 口径检查接口、搜索、首屏、匹配、通知和并发认领指标。
  \item \textbf{接口一致性：}确认参数命名、响应结构、错误码和枚举与 API 文档及前端类型一致。
  \item \textbf{可维护质量：}通过格式、静态类型、单元测试和构建检查，保证变更可被稳定集成和回归。
\end{enumerate}

\subsection{测试策略与层次}

项目采用由静态检查、单元测试、服务测试、接口测试和场景验收组成的分层策略。靠近代码的测试负责快速定位规则错误，靠近用户场景的测试负责检查跨模块业务链路；高频核心流程、权限规则和状态转换优先自动化。

\begin{figure}[H]
  \centering
  \begin{tikzpicture}[node distance=3.5mm]
    \node[draw=reportgray,fill=gray!7,rounded corners=2pt,minimum width=14.2cm,minimum height=0.9cm] (static)
      {静态与构建检查：Ruff、Mypy、vue-tsc、前端构建与镜像构建};
    \node[draw=sysugreen,fill=sysugreenlight,rounded corners=2pt,minimum width=12.6cm,minimum height=0.9cm,above=of static] (unit)
      {单元测试：纯函数、校验规则、状态转换、评分与格式化};
    \node[draw=reportblue,fill=blue!5,rounded corners=2pt,minimum width=10.8cm,minimum height=0.9cm,above=of unit] (service)
      {服务与仓储测试：事务协作、权限、幂等和异常分支};
    \node[draw=reportorange,fill=orange!7,rounded corners=2pt,minimum width=9.0cm,minimum height=0.9cm,above=of service] (api)
      {接口契约测试：鉴权、参数、响应和错误码};
    \node[draw=reportred,fill=red!5,rounded corners=2pt,minimum width=7.2cm,minimum height=0.9cm,above=of api] (scenario)
      {场景验收：多角色核心业务流程};
    \draw[flow] ([xshift=7.5mm]static.east) -- node[right,font=\zihao{-5},align=left]{业务完整性提高\\定位范围扩大} ([xshift=7.5mm]scenario.east);
  \end{tikzpicture}
  \caption{测试层次与职责}
\end{figure}

各层采用的主要方法如下：

\begin{table}[H]
  \centering
  \caption{测试类型与实施方法}
  \begin{tabularx}{\textwidth}{L{2.5cm}L{3.5cm}Y}
    \toprule
    类型 & 实施方式 & 适用内容 \\
    \midrule
    静态检查 & Ruff、Mypy、vue-tsc & Python 风格和类型、异步调用、TypeScript 类型与导入关系 \\
    单元测试 & Pytest、Vitest & 业务规则、工具函数、评分计算、配置校验和错误分类 \\
    服务测试 & 依赖替身、测试数据库 & Service 与 Repository 协作、状态机、事务、通知和对象引用 \\
    接口测试 & FastAPI TestClient/ASGI Client & 路由鉴权、参数边界、状态码、响应 Schema 和角色权限 \\
    场景测试 & 按角色执行核心用例 & 注册、发布、匹配、认领、交接、审核等端到端业务步骤 \\
    性能测试 & 固定数据集、并发请求与受控浏览器计时 & 接口 P95、搜索容量、首屏、匹配、通知和并发认领 \\
    回归测试 & 缺陷对应测试加相邻路径测试 & 修复本身、同类输入、上下游状态与公共模块调用方 \\
    \bottomrule
  \end{tabularx}
\end{table}

\subsection{测试设计方法}

\begin{enumerate}
  \item \textbf{等价类划分：}将认证状态、角色、物品类型和业务状态划分为有效类与无效类，各选代表输入。
  \item \textbf{边界值分析：}检查验证码长度、分页参数、时间区间、匹配分数、上传数量和文本长度的边界。
  \item \textbf{判定表：}组合登录状态、资源归属、角色权限和目标状态，检查允许、拒绝及冲突返回。
  \item \textbf{状态迁移：}依据认领与交接状态图覆盖合法迁移、重复操作、逆向迁移和终止状态。
  \item \textbf{场景法：}按用户目标串联跨模块步骤，验证发布至找回、举报至处置等完整业务链路。
  \item \textbf{异常推测：}针对网络失败、AI 不可用、重复点击、过期数据和外部服务异常设计用例。
\end{enumerate}

\subsection{环境与测试数据}

\begin{table}[H]
  \centering
  \caption{测试环境配置}
  \begin{tabularx}{\textwidth}{L{3.1cm}Y}
    \toprule
    环境项 & 配置说明 \\
    \midrule
    后端运行环境 & Python 3.12、FastAPI、SQLAlchemy；单元与服务测试使用隔离测试数据库 \\
    AI 服务环境 & Python 3.12、FastAPI；外部模型调用使用可控替身并检查异常分支 \\
    前端运行环境 & Node.js、pnpm、Vue 3、Vite、Vitest、vue-tsc \\
    集成配置 & MySQL、MinIO、Backend、AI 服务、用户端和管理端使用 Compose 统一描述 \\
    性能基准 & 50 名并发活跃用户、10 万条物品记录、每页 20 条；页面计时使用 20 Mbps 受控网络 \\
    数据隔离 & 每个测试独立准备账号、物品与状态，执行后回滚或清理，禁止复用个人真实资料 \\
    \bottomrule
  \end{tabularx}
\end{table}

测试数据至少包含两个普通用户、一个管理员、已认证与未认证账号、普通物品与证件类物品，以及处于待审核、匹配中、认领中、已交接和已关闭状态的记录。边界数据覆盖空字符串、最大长度、非法枚举、反向时间区间、重复请求和无权资源标识。敏感图片使用专用样例，不使用真实证件或个人隐私资料。

\subsection{阶段安排与人员职责}

\begin{table}[H]
  \centering
  \caption{测试阶段安排}
  \begin{tabularx}{\textwidth}{L{2.2cm}L{3.4cm}Y L{3.1cm}}
    \toprule
    阶段 & 输入 & 主要活动 & 输出 \\
    \midrule
    测试准备 & 需求基线、模型、接口文档 & 划分范围与优先级，确定环境、角色、数据和准入条件 & 测试计划、用例清单 \\
    模块测试 & 可独立运行的模块 & 编写单元、服务和接口测试，执行静态检查 & 测试记录、缺陷单 \\
    联调测试 & 前后端及 AI 契约 & 检查参数、枚举、错误码和跨模块状态传递 & 联调问题记录 \\
    回归测试 & 修复版本与缺陷单 & 复测原场景并检查相邻流程，执行自动化测试集 & 回归结果、关闭记录 \\
    版本验收 & 稳定基线与核心数据 & 按 18 个核心用例和 6 个性能用例检查业务、权限与质量指标 & 验收记录、报告 \\
    \bottomrule
  \end{tabularx}
\end{table}

测试统筹由朱梓涵与程文韬共同负责，包括计划维护、用例整理、缺陷汇总和版本验收；周星辰负责用户端与管理端测试；朱梓涵负责用户、认证、物品和集成测试；郭焜华负责认领、交接、通知、积分和治理测试；林洪宽负责 AI 服务与匹配测试。缺陷修复由对应模块负责人实施，并由另一名成员复核回归结果。

\subsection{进入、退出与暂停准则}

\begin{table}[H]
  \centering
  \caption{测试控制准则}
  \begin{tabularx}{\textwidth}{L{2.5cm}Y}
    \toprule
    准则 & 内容 \\
    \midrule
    进入条件 & 需求和接口形成基线；测试模块可安装运行；数据库迁移可执行；测试账号与数据可准备 \\
    退出条件 & 核心用例全部通过；阻断级与严重级缺陷关闭；相关静态检查、自动化测试和构建通过 \\
    暂停条件 & 环境不可用、迁移失败、测试数据污染，或出现可能影响其他用例判断的阻断级缺陷 \\
    恢复条件 & 环境恢复到隔离基线；根因已定位；受影响用例能够重复执行并得到稳定结果 \\
    \bottomrule
  \end{tabularx}
\end{table}

\section{核心测试用例}

表 \ref{tab:core-cases} 给出课程项目的 18 个功能、安全与一致性核心用例。用例既覆盖正常业务主线，也覆盖越权、重复操作、状态冲突和敏感资源访问。

\zihao{-5}
\begin{longtable}{L{2.5cm}L{3.6cm}L{3.9cm}L{3.9cm}}
  \caption{核心测试用例}\label{tab:core-cases}\\
  \toprule
  编号与目标 & 前置条件与操作 & 预期结果 & 执行结果 \\
  \midrule
  \endfirsthead
  \multicolumn{4}{c}{\tablename\ \thetable（续）}\\
  \toprule
  编号与目标 & 前置条件与操作 & 预期结果 & 执行结果 \\
  \midrule
  \endhead
  \midrule
  \multicolumn{4}{r}{续下页}\\
  \endfoot
  \bottomrule
  \endlastfoot

  \shortstack[l]{TC-LOGIN-01\\用户注册登录} & 使用未注册手机号获取验证码，填写昵称和密码注册，再使用密码登录 & 创建普通用户并返回有效会话；验证码一次性消费；错误、过期及重复验证码被拒绝 & 注册、密码登录、验证码校验与会话加载均通过 \\
  \shortstack[l]{TC-LOGIN-02\\管理员登录} & 使用管理员账号登录后台，再用普通用户访问管理接口 & 管理员进入后台；普通用户请求被拒绝；管理员权限由后端判定 & 管理员鉴权与普通用户越权检查通过 \\
  \shortstack[l]{TC-CERT-01\\提交校园认证} & 未认证用户填写姓名、校园编号并提交认证图片 & 生成待审核申请；图片归属正确；重复提交按状态规则处理 & 申请创建、字段校验和资源归属检查通过 \\
  \shortstack[l]{TC-CERT-02\\认证审核} & 管理员查看待审核申请并选择通过或驳回 & 状态只从待审核迁移一次；意见被保存；用户收到结果通知 & 通过、驳回、重复审核和通知检查通过 \\
  \shortstack[l]{TC-LOST-01\\发布失物} & 登录用户填写名称、类别、时间、地点、描述和图片并提交 & 生成失物记录；必填项和时间区间合法；发布者可查看详情 & 发布、校验、上传引用和详情查询通过 \\
  \shortstack[l]{TC-FOUND-01\\发布招领} & 登录用户填写招领信息、保管方式、验证问题和图片 & 生成招领记录；验证问题完整；敏感图片按访问规则处理 & 发布、问题保存、审核状态和图片保护检查通过 \\
  \shortstack[l]{TC-SEARCH-01\\搜索筛选} & 准备不同类别、地点、时间和状态的物品，组合筛选并翻页 & 返回结果满足全部条件；排序和分页稳定；不可见记录被过滤 & 类别、状态、地点、时间、排序和分页检查通过 \\
  \shortstack[l]{TC-MATCH-01\\生成匹配} & 准备发布者不同、内容相近且状态有效的失物和招领 & 形成唯一匹配对；分数位于合法范围；重复计算不产生重复有效记录 & 评分、去重、状态和通知检查通过 \\
  \shortstack[l]{TC-CLAIM-01\\认领验证通过} & 失主从关联匹配发起认领，以同义、口语或轻微错字回答全部验证问题 & 校验匹配归属与答案完整性；小模型语义分通过后进入合法认领状态；不可用时标记规则降级 & 归属、语义答案、来源记录、审核与状态迁移检查通过 \\
  \shortstack[l]{TC-CLAIM-02\\认领被拒绝} & 拾获者查看进行中的认领，填写原因并拒绝 & 认领转为已拒绝；招领和匹配按规则释放；认领人收到通知 & 拒绝、资源释放和通知检查通过 \\
  \shortstack[l]{TC-CLAIM-03\\重复认领} & 对同一招领连续或并发提交多个有效认领请求 & 系统只保留一个有效认领；其余请求得到明确冲突响应 & 状态检查、条件更新和唯一性约束检查通过 \\
  \shortstack[l]{TC-HO-01\\双方交接} & 已批准认领已建立交接安排，双方分别确认并重复点击 & 确认操作幂等；双方确认后认领、失物、招领和匹配同步结案 & 双方确认、重复确认和结案状态检查通过 \\
  \shortstack[l]{TC-CREDIT-01\\积分流水} & 完成交接并再次提交相同完成请求 & 双方积分按规则变化；流水记录原因与实际变化；重复请求不重复记分 & 积分更新、流水和幂等检查通过 \\
  \shortstack[l]{TC-ADMIN-01\\治理审计} & 管理员查看内容与全部状态认领，打开问答详情后执行合法审核 & 详情展示参考答案快照、认领回答、分数、来源、凭证和交接；普通用户响应不含参考答案 & 管理员详情、隐私隔离、审核日志和重复操作检查通过 \\
  \shortstack[l]{TC-REPORT-01\\举报处理} & 用户选择明确目标提交举报，管理员核实并处置 & 举报状态更新；目标内容与账号按处置规则变化；相关用户收到通知 & 举报创建、目标关联、处置和通知检查通过 \\
  \shortstack[l]{TC-SEC-01\\未登录访问} & 不携带会话访问个人资料、发布、认领和管理接口 & 请求被拒绝；业务数据不变化；响应不包含敏感内容 & 各类受保护路由鉴权检查通过 \\
  \shortstack[l]{TC-SEC-02\\普通用户越权} & 普通用户访问他人敏感资源或调用审核、治理接口 & 服务端拒绝请求；仅隐藏前端按钮不能绕过权限 & 资源归属和角色权限检查通过 \\
  \shortstack[l]{TC-SEC-03\\敏感图片访问} & 匿名用户、无关用户和有权角色分别请求证件与认领凭证 & 仅有权角色取得短期访问地址；无权请求不返回原图 & 响应字段、资源归属和签发权限检查通过 \\
\end{longtable}
\normalsize

\subsection{性能测试用例}

性能测试沿用需求分析报告的负载、数据规模和计时边界。每个场景先完成预热，再在稳定时段重复采样；服务端场景统计请求或任务的 P95，页面场景统计关键内容可见时间，并记录错误率、资源占用和异常响应。

\small
\begin{longtable}{L{2.1cm}L{4.2cm}L{4.5cm}L{3.3cm}}
  \caption{性能需求对应测试用例}\label{tab:performance-cases}\\
  \toprule
  编号与需求 & 场景与负载 & 操作与采样方法 & 通过条件 \\
  \midrule
  \endfirsthead
  \multicolumn{4}{c}{\tablename\ \thetable（续）}\\
  \toprule
  编号与需求 & 场景与负载 & 操作与采样方法 & 通过条件 \\
  \midrule
  \endhead
  \bottomrule
  \endlastfoot
  \shortstack[l]{PT-01\\PERF-01} & 50 名活跃用户混合执行查询、创建和状态操作，不含上传与 AI 计算 & 预热后持续采样普通业务接口，统计响应时间和非预期错误 & P95 不超过 500 ms；错误率不高于 1\% \\
  \shortstack[l]{PT-02\\PERF-02} & 10 万条物品记录，每页 20 条，组合类别与地点筛选并按时间排序 & 轮换关键词、筛选组合和页码，统计搜索请求的服务端耗时 & P95 不超过 1 s；分页结果稳定 \\
  \shortstack[l]{PT-03\\PERF-03} & 20 Mbps 网络和主流移动设备视口，访问首页、搜索页和详情页 & 清理页面缓存后分别导航 30 次，记录关键内容首次可见时间 & 95\% 的页面在 1.5 s 内显示关键内容 \\
  \shortstack[l]{PT-04\\PERF-04} & AI 服务可用，同时提交不超过 20 个候选匹配任务 & 记录任务接受时间与结果可查询时间，覆盖文本、地点、时间和图片组合 & 95\% 的匹配在 3 s 内形成结果 \\
  \shortstack[l]{PT-05\\PERF-05} & 连续触发匹配、认领、审核、交接和积分事件 & 记录业务状态写入与通知、未读数可查询的时间差 & 95\% 的通知在 2 s 内可查询 \\
  \shortstack[l]{PT-06\\PERF-06} & 10 个用户同时认领同一招领，重复执行 20 轮 & 同步释放并发请求，检查响应时间、活动认领数量和状态一致性 & 请求在 2 s 内返回；每轮最多一个活动认领 \\
\end{longtable}
\normalsize

\subsection{需求与用例对应关系}

\begin{table}[H]
  \centering
  \caption{需求域与测试用例对应关系}
  \begin{tabularx}{\textwidth}{L{3.2cm}L{5.1cm}Y}
    \toprule
    需求域 & 对应用例 & 主要检查点 \\
    \midrule
    账号与认证 & TC-LOGIN-01/02、TC-CERT-01/02 & 登录方式、角色、申请审核、资源归属 \\
    物品发布与搜索 & TC-LOST-01、TC-FOUND-01、TC-SEARCH-01 & 字段校验、状态、上传、筛选和分页 \\
    匹配与认领 & TC-MATCH-01、TC-CLAIM-01/02/03 & 评分、去重、归属、问答和冲突 \\
    交接与激励 & TC-HO-01、TC-CREDIT-01 & 状态结案、幂等、积分和流水 \\
    平台治理 & TC-ADMIN-01、TC-REPORT-01 & 审核、举报、处置、日志和通知 \\
    安全与隐私 & TC-SEC-01/02/03 & 认证、授权、敏感图片与服务端边界 \\
    性能与容量 & PT-01 至 PT-06 & P95、数据规模、首屏、匹配、通知和并发一致性 \\
    \bottomrule
  \end{tabularx}
\end{table}

\section{测试执行与结果}

\subsection{执行流程}

每次测试从固定基线开始：安装锁定依赖，初始化测试数据库，执行静态检查和模块测试，再运行跨模块用例与构建检查。测试失败时记录环境、输入、步骤、实际结果和日志摘要；修复合入后先复测原用例，再执行同模块及相邻业务路径。测试数据在用例结束后清理，避免前一用例改变后一用例的判定。

\subsection{自动化结果摘要}

期末基线的自动化测试按四个代码模块执行，结果汇总见表 \ref{tab:test-results}。这些数字用于说明测试执行范围，不替代核心用例中的业务判定。

\begin{table}[H]
  \centering
  \caption{自动化测试结果摘要}\label{tab:test-results}
  \begin{tabularx}{\textwidth}{L{2.8cm}C{2.2cm}L{3.5cm}Y}
    \toprule
    模块 & 用例数 & 结果 & 主要覆盖 \\
    \midrule
    主后端 & 204 & 全部通过，语句覆盖率 80\% & 用户、物品、匹配、语义认领、交接、积分、通知、治理和迁移 \\
    AI 服务 & 64 & 全部通过 & 配置、鉴权、分类、敏感检测、匹配评分、回答语义校验和异常处理 \\
    用户端 & 17 & 全部通过 & 会话、导航、筛选参数、上传辅助和错误分类 \\
    管理端 & 15 & 全部通过 & 管理会话、列表请求、举报参数、错误分类和格式化 \\
    \midrule
    合计 & 300 & 全部通过 & 四个代码模块的自动化回归集合 \\
    \bottomrule
  \end{tabularx}
\end{table}

静态检查同时覆盖 Python 格式与类型、TypeScript 类型和前端生产构建。测试结果与 18 个核心用例共同用于版本退出判定：自动化集合负责稳定回归，核心用例负责确认业务语义和角色边界。

\section{缺陷跟踪}

\subsection{缺陷记录字段与分级}

每条缺陷统一记录编号、标题、发现版本、环境、前置条件、复现步骤、预期结果、实际结果、严重度、负责人、修复提交、回归用例和当前状态。相同根因导致的多个现象合并跟踪，不同业务状态或权限边界分别建立条目。

\begin{table}[H]
  \centering
  \caption{缺陷严重度定义}
  \begin{tabularx}{\textwidth}{L{1.8cm}L{3.1cm}Y L{3.0cm}}
    \toprule
    等级 & 名称 & 判定标准 & 处理规则 \\
    \midrule
    S1 & 阻断级 & 主流程无法继续，或出现越权、隐私泄露、严重数据矛盾 & 暂停版本验收，立即修复并执行专项回归 \\
    S2 & 严重级 & 关键功能异常、常见操作断链或重要契约不一致 & 当前迭代修复，补充回归用例后关闭 \\
    S3 & 一般级 & 次要功能、低频边界或局部交互异常 & 排入迭代，按影响范围执行回归 \\
    S4 & 提示级 & 文案、样式或不影响功能的易用性问题 & 记录后集中处理并复查界面 \\
    \bottomrule
  \end{tabularx}
\end{table}

严重度表示缺陷造成的影响，优先级表示处理顺序。默认情况下 S1 对应最高优先级；若缺陷影响课程演示主线或多个模块，即使严重度相同也优先处理。

\subsection{缺陷闭环流程}

\begin{figure}[H]
  \centering
  \begin{tikzpicture}[node distance=10mm and 12mm]
    \node[state] (discover) {发现与复现\\记录环境/输入};
    \node[state,right=of discover] (classify) {确认与分级\\指定负责人};
    \node[state,right=of classify] (fix) {定位与修复\\代码评审};
    \node[state,below=of fix] (regress) {原用例复测\\相邻路径回归};
    \node[state,left=of regress] (gate) {质量门禁\\静态/测试/构建};
    \node[state,left=of gate] (close) {验收关闭\\保留记录};
    \draw[flow] (discover) -- (classify);
    \draw[flow] (classify) -- (fix);
    \draw[flow] (fix) -- (regress);
    \draw[flow] (regress) -- (gate);
    \draw[flow] (gate) -- (close);
    \draw[flow] (close.west) -- ++(-0.65,0) |- node[left,font=\zihao{-5}]{再次出现} (discover.west);
    \draw[dashedflow] (gate.north) -- node[right,font=\zihao{-5}]{未通过} (fix.south);
  \end{tikzpicture}
  \caption{缺陷跟踪与回归闭环}
\end{figure}

缺陷仅在以下条件同时满足时关闭：复现步骤不再出现异常；修复通过代码评审；至少新增或更新一个回归用例；相邻状态与公共调用方未受影响；静态检查、自动化测试和构建均通过。再次出现同类现象时重新打开原条目，并补充新的触发条件。

\subsection{代表性缺陷记录}

\zihao{-5}
\begin{longtable}{L{1.3cm}L{3.0cm}L{3.0cm}L{4.0cm}L{2.6cm}}
  \caption{代表性缺陷闭环记录}\label{tab:defects}\\
  \toprule
  编号 & 缺陷现象 & 根因 & 修复与回归 & 状态 \\
  \midrule
  \endfirsthead
  \multicolumn{5}{c}{\tablename\ \thetable（续）}\\
  \toprule
  编号 & 缺陷现象 & 根因 & 修复与回归 & 状态 \\
  \midrule
  \endhead
  \bottomrule
  \endlastfoot
  D-01 & 验证码可重复使用，调试码暴露范围过大 & 验证码生命周期与环境开关约束不足 & 改为随机、限时、一次性消费；增加错误、过期、重复消费和白名单用例 & 已关闭 \\
  D-02 & 无权用户可能取得敏感图片地址 & 业务记录直接暴露对象地址，缺少统一授权签发 & 保存对象引用并按角色签发短期地址；增加匿名、无关用户和有权角色对照用例 & 已关闭 \\
  D-03 & 同一招领可能产生多个有效认领 & 先查询后创建缺少状态条件和唯一性保护 & 增加条件更新、状态检查和数据库唯一约束；加入重复与并发请求回归 & 已关闭 \\
  D-04 & 双方确认后部分关联状态未同步结案 & 交接完成条件与多实体迁移分散在不同分支 & 将确认设计为幂等操作并统一结案事务；回归双方顺序与重复确认 & 已关闭 \\
  D-05 & 发布物品后未自动触发匹配计算 & 发布事务与异步任务调度没有形成连接 & 在同一事务登记持久任务并由执行器处理；回归失物与招领两种发布路径 & 已关闭 \\
  D-06 & 前后端对通知已读响应的结构理解不同 & API Schema 与 TypeScript 类型未同步 & 统一响应结构与共享类型；增加路由契约和前端解析测试 & 已关闭 \\
  D-07 & 同义或口语化认领答案因关键词表述不同被拒绝 & 认领验证只做精确关键词包含判断 & 增加小模型语义评分和关键词降级；覆盖同义表达、轻微错字与明确否定 & 已关闭 \\
  D-08 & 管理后台只能查看申诉，无法还原完整问答证据 & 管理员列表默认过滤申诉，回答记录未保存参考答案与来源快照 & 改为全状态认领治理；保存并仅向管理员展示问答快照、分数和来源 & 已关闭 \\
\end{longtable}
\normalsize

\section{质量保证方法}

\subsection{评审与追踪}

\begin{enumerate}
  \item \textbf{需求评审：}检查需求是否具有唯一编号、明确角色、输入条件、正常流程、异常流程和可判定结果。
  \item \textbf{设计评审：}用用例图、类图、序列图和状态图检查模块边界、对象关系与状态迁移是否一致。
  \item \textbf{契约评审：}API、枚举、数据库字段和前端类型同步检查；接口变化同时更新调用方与测试。
  \item \textbf{代码评审：}重点检查权限入口、事务边界、异常处理、外部调用、迁移脚本和公共组件影响范围。
  \item \textbf{测试追踪：}功能域对应核心用例，缺陷对应回归用例，使需求、测试、修改和版本结果能够相互定位。
\end{enumerate}

\subsection{自动化质量检查}

\begin{table}[H]
  \centering
  \caption{质量检查设置}
  \begin{tabularx}{\textwidth}{L{2.7cm}L{4.3cm}Y}
    \toprule
    检查项 & 检查内容 & 判定与用途 \\
    \midrule
    Python 静态检查 & Ruff 规则、格式、Mypy 严格类型 & 无 lint、格式和类型错误 \\
    前端静态检查 & vue-tsc、模块导入和共享类型 & 两个前端均通过类型检查 \\
    自动化测试 & Backend、AI 服务、用户端、管理端测试集 & 全部用例通过，无失败与异常退出 \\
    覆盖率报告 & Backend 语句覆盖率与核心模块分布 & 生成覆盖率报告，对照 80\% 目标复核关键规则 \\
    构建检查 & 两个前端生产构建、Backend 与 AI 镜像构建 & 构建命令成功，产物可生成 \\
    配置检查 & Compose 展开与环境变量映射 & 配置可解析，必填变量具有明确校验 \\
    \bottomrule
  \end{tabularx}
\end{table}

\subsection{安全与隐私检查}

安全检查以服务端边界为准，覆盖密码哈希、验证码生命周期、管理员初始化、内部服务鉴权、CORS 白名单、URL 获取限制、对象所有权、私有存储和短期访问地址。敏感图片默认受保护，只有明确有权的角色才能取得访问地址；AI 返回先经过结构与范围校验，再进入业务判断。

安全相关缺陷修复后执行专项回归，并同时检查登录、认证、物品、认领和管理接口，避免公共鉴权或对象访问逻辑改变其他角色行为。测试日志和错误响应不记录密码、验证码、令牌及图片原始内容。

\subsection{一致性与回归保证}

认领、交接、积分、通知和匹配属于共享状态链路。测试同时覆盖合法迁移、非法迁移、重复请求和并发冲突，检查条件更新、唯一约束、事务提交和幂等键是否共同生效。数据库结构变更使用 Alembic 管理，并测试迁移顺序、约束建立和单一 head。

回归范围按影响传播确定：局部工具函数修改执行所在模块测试；接口或共享类型修改同时执行前后端相关测试；状态机、权限、迁移和对象存储修改执行主后端全量测试及对应场景用例。每条关闭缺陷均保留固定回归用例，使同类问题能够在以后版本中自动被发现。

\section{结论}

本报告建立了覆盖五类测试对象、六种测试类型、18 个功能与安全核心场景及 6 个性能场景的测试方案，明确了测试环境、数据、阶段、人员、准入退出和执行方法。期末基线的四个模块共 300 个自动化用例全部通过，Backend 语句覆盖率达到 80\%，核心业务、权限和状态场景按统一预期结果完成检查。

项目同时采用统一缺陷字段、四级严重度、修复回归闭环和自动化质量门禁。测试计划、测试用例、缺陷跟踪与质量保证方法共同构成课程项目的质量控制体系，并满足本报告规定的测试目标。
